\subsubsection{止まる値を調べる}\label{節：止まる値を調べる}
次の値が今の値と等しい\bookterm{fixed}{不動点}に達すると,\bookterm{sequence}{数列}は止まる.
丸めを外した実数の\bookterm{recurrence}{漸化式}とは,止まる値が異なることがある.

整数$K$を引いて,
\[
  b_{n+1}=\left\lfloor\frac{b_n}{2}\right\rfloor
\]
に帰着できる場合がある.
正整数$m$に対して,
\[
  \left\lfloor\frac{\lfloor x\rfloor}{m}\right\rfloor
  =\left\lfloor\frac xm\right\rfloor
\]
が成り立つ.$\lfloor x\rfloor=mq+r\ (0\le r<m)$と書けば,$x=mq+r+t\ (0\le t<1)$.
$0\le r+t<m$より,両方の\bookterm{floor-function}{床関数}が$q$となるからである.

\textbf{負数にも注意する.}$-1<x<0$なら$\lfloor x\rfloor=-1$である.

\noindent \bookblocklink{example-gauss-role-6}{problem}{answer}{【例題】}\par\nobreak
\begin{enumerate}[label=(\arabic*),parsep=3pt,beginpenalty=10000,format=\bookitemlink{example-gauss-role-6}{problem}{answer}]
  \item 非負整数$M$に対し,\bookterm{sequence}{数列}を次のように定める.
    \[
      a_1=M,\qquad
      a_{n+1}=\left\lfloor\frac{a_n+6}{2}\right\rfloor
      \quad(n\ge1).
    \]
    \bookterm{general}{一般項}と,最終的に一定になる値を求めよ.
\end{enumerate}

\noindent \bookblocklink{example-gauss-role-6}{hint}{problem}{【方針】}\par\nobreak
\begin{enumerate}[label=(\arabic*),parsep=3pt,beginpenalty=10000,format=\bookitemlink{example-gauss-role-6}{hint}{problem}]
  \item 6を引くと,\bookterm{floor-function}{床関数}によって半分にする\bookterm{recurrence}{漸化式}になります.
    初項が6より小さい場合は,負の数の\bookterm{floor-function}{床関数}を使って考えましょう.
\end{enumerate}

\noindent \bookblocklink{example-gauss-role-6}{answer}{problem}{【解答】}\par\nobreak
\begin{enumerate}[label=(\arabic*),parsep=3pt,beginpenalty=10000,format=\bookitemlink{example-gauss-role-6}{answer}{problem}]
  \item \[
      a_n=6+\left\lfloor\frac{M-6}{2^{n-1}}\right\rfloor.
    \]
    有限回の操作の後で,$M<6$なら5,$M\ge6$なら6に固定される.
\end{enumerate}

\noindent \bookblocklink{example-gauss-role-6}{solution}{problem}{【解法】}\par\nobreak
\begin{enumerate}[label=(\arabic*),parsep=3pt,beginpenalty=10000,format=\bookitemlink{example-gauss-role-6}{solution}{problem}]
  \item 初項からすべての項が非負整数となる.
    $b_n=a_n-6$とおくと,整数6を\bookterm{floor-function}{床関数}の外に出せるので,
    \begin{align*}
      b_{n+1}
      &=\left\lfloor\frac{b_n+12}{2}\right\rfloor-6\\
      &=\left\lfloor\frac{b_n}{2}\right\rfloor.
    \end{align*}
    \bookterm{floor-function}{床関数}の合成の性質と帰納法により,
    \[
      b_n=\left\lfloor\frac{M-6}{2^{n-1}}\right\rfloor.
    \]
    十分大きな$n$で,中身は$M\ge6$なら$[0,1)$,$M<6$なら$(-1,0)$に入る.
    床はそれぞれ0,$-1$なので,6または5に有限回で到達する.

    実際,整数$x$が\bookterm{fixed}{不動点}となる条件は,
    \[
      x\le\frac{x+6}{2}<x+1
      \iff 4<x\le6
    \]
    より5と6が\bookterm{fixed}{不動点}となる.初項による分岐は\sectionref{節：初項による違い}でも見る.
\end{enumerate}
